\chapter{AR Pathfinder}
\label{ch:ar-pathfinder}
% Port from: docs/10_ar_pathfinder.md

The pathfinder combines four real-time data streams --- GPS coordinates,
device compass heading, WebSocket proximity messages, and camera-based AR
rendering --- into one navigation experience. Code listings for the formulas
referenced here appear in Appendix~\ref{app:code}.

\section{Geofencing and Two-Phase Arrival}
All proximity checks run server-side using the haversine great-circle
formula, accurate to well under a metre at campus distances --- far finer
than the 3--5\,m accuracy of consumer GPS itself. A waypoint counts as
reached when the reported position falls within the place's geofence
radius. The radius --- the \code{visit\_threshold\_meters} field on the
place --- defaults to 10\,m rather than being hard-coded: dense
environments may need a
tighter fence to avoid false arrivals, open spaces a looser one to absorb
GPS drift, and administrators can tune it without a deployment.

Arrival is deliberately a two-phase process. Detection --- the server
reporting \code{arrived: true} --- only raises a confirmation prompt.
The checkpoint is marked visited only when the user explicitly confirms,
and the server re-validates the last reported position against the same
radius before accepting. A momentary GPS glitch therefore cannot complete a
checkpoint, and a confirmed visit is backed by two independent
within-radius checks.

\section{Compass Guidance}
The server includes a true-north bearing to the next waypoint in every
progress frame. The client subtracts the device's compass heading to obtain
the relative bearing that drives the HUD arrow. Two engineering details
make this usable. First, raw \code{DeviceOrientationEvent} readings jitter
by several degrees per frame, so headings pass through an exponential
moving average (weight 0.15 on the new reading) that removes jitter while
still responding to turns within about two seconds. Second, orientation
events arrive at up to 60\,Hz; routing them through React state would
re-render the component tree at that rate, so the arrow and distance text
are updated imperatively through refs, bypassing rendering entirely. iOS
additionally requires an explicit permission request for orientation data,
which is why navigation can only start from a button press (a user
gesture).

\section{AR Camera Layer}
The AR view uses A-Frame with AR.js location-based tracking, with the
camera feed rendered as a WebGL texture. The \code{<a-scene>} element is
injected directly under \code{<body>}, outside the React tree --- A-Frame
requires this to size its fullscreen canvas correctly --- and is fully torn
down (including stopping webcam tracks) when navigation ends. The vendor
scripts are served locally and loaded on demand only when AR mode starts,
which pins versions (A-Frame 1.6.0, AR.js 3.4.7), avoids ad-blocker
interference, and keeps the initial page load small.

Milestone and side-quest videos play on a plane anchored in world space.
Anchoring by GPS coordinates would make the plane jump with satellite
drift, so instead the plane spawns four metres in front of the camera, its
rotation frozen at spawn (looking around pans off it, as with a real
object), while its position is re-glued to the camera's translation every
animation frame. The result reads as ``placed in the world'' yet is immune
to GPS jitter.

\section{Waypoint Types}
\begin{table}[h]
  \centering\small
  \begin{tabular}{@{}lp{0.68\textwidth}@{}}
    \toprule
    \textbf{Type} & \textbf{Behaviour on arrival} \\
    \midrule
    \code{start} & Confirmation prompt; opening image shown \\
    \code{node} & Turn point; directional image shown on confirmation \\
    \code{milestone} & Video plays; confirmation awards points and logs the milestone \\
    \code{side\_quest} & Video plays; accept starts the 5-minute dwell timer, skip continues \\
    \code{end} & Final checkpoint; completion image and trip-complete state \\
    \bottomrule
  \end{tabular}
  \caption{Waypoint type semantics.}
  \label{tab:waypoints}
\end{table}

\section{Graceful Degradation and Demo Mode}
The design principle is that the compass arrow is always the fallback.
Without HTTPS (camera APIs require a secure context) or if the AR scripts
fail to load, navigation continues with the compass HUD and a visible
notice; if video autoplay is blocked, a tap-to-play overlay appears; if the
compass reports non-absolute headings, a calibration banner is shown.
A separate demo mode bypasses GPS entirely: it replays the real route from
hardcoded waypoints with a ``simulate arrival'' control, exercising the
complete flow --- media, confirmations, rewards UI --- indoors, which is
essential for stage demonstrations and QA.

\section{The Seeded Route}
The production route walks Orchid College over approximately 400\,m
(five to seven minutes): a start point, three turn nodes, two milestones
with video stories (the parking area and Orchid International College), one
side quest (Bishwajit Medical Hall), and an end point, with a 10\,m
geofence per waypoint and the college badge on completion.
